\section{Uniform cohomology for tame coefficient systems}\label{sec:tame-coefficients}

We prove the coefficient estimate by bounding poles on a resolution. The problem is that tame inertia can have arbitrarily large order. Passing to a root stack makes the sheaf unramified; a character count will remove the dependence on the chosen root order.

Keep the notation of Theorem~\ref{thm:local-coefficients}. Write $B_{ij}=C_i\cdot C_j$, $g_i=g(C_i)$, and let $\nu_i$ be the number of crossing points on $C_i$. The matrix $B$ is negative definite \cite[Lemma 14.1, p.~224]{Lipman} and has nonnegative off-diagonal entries. Choose positive integers $c_i$ such that
\[
 Z=\sum_i c_iC_i,\qquad J=\OO_X(-Z),\qquad
 d_i=\deg(J|_{C_i})=-\sum_jc_jB_{ij}>0.
\]
Such integers exist because $(-B)^{-1}$ is entrywise nonnegative, as proved in Section~\ref{sec:abelian-inertia}; its product with the all-ones vector is strictly positive. Clear denominators. Put
\begin{equation}\label{eq:resolution-constants}
 L_i=\sum_jc_j|B_{ij}|,\qquad
 K_i=\max\{0,-B_{ii},2g_i-2+\nu_i\},\qquad
 b_i=\frac{2(p-1)K_i}{p}.
\end{equation}
We will use an integer $m_0\geq2$ larger than every $(b_i+L_i)/d_i$.

\subsection{Removing tame inertia}

Choose $N$ prime to $p$, divisible by all inertia orders of the monodromy of $M$, and form the stack $\mathfrak X$ obtained by taking $N$th roots of each $C_i$. Write $\pi:\mathfrak X\to X$ and $\mathfrak C_i$ for the reduced divisors above $C_i$. At a crossing with parameters $x,y$, its chart is
\[
 \left[\Spec R[u,v]/(u^N-x,v^N-y)\,/\,\mu_N^2\right].
\]
Away from crossings there is one root variable. These regular charts show that $\mathfrak X$ is a regular tame Deligne--Mumford stack, proper over $A$, with coarse space $X$ and $N\mathfrak C_i=\pi^*C_i$. This is the usual root construction \cite[\S\S2.3--2.4, author pp.~7--10]{Cadman}.

The sheaf $M$ extends to a locally constant sheaf on $\mathfrak X$. The root charts kill its tame codimension-one inertia by Abhyankar's lemma \cite[Expos\'e XIII, Proposition 5.2, original pp.~428--429]{SGA1}. Normalize each regular chart in the corresponding finite cover. Purity of the branch locus makes the normalization \'etale, and uniqueness of normalization gives the descent data on overlaps \cite[Tags 0BRM and 0BMB]{Stacks}. This argument places no prime-to-$p$ restriction on the entire monodromy group.

Set $V=M\otimes_{\F_p}\OO_{\mathfrak X}$ and write $r=\operatorname{rank}M$. It is a vector bundle with $F^*V\simeq V$ for absolute Frobenius, and is trivialized by a finite \'etale cover. The divisor $\mathfrak C_i$ is a $\mu_N$-gerbe over the curve obtained from $C_i$ by taking roots at its crossing points. Normalize degree on it so that pullback from $C_i$ preserves ordinary degree. A line bundle has its $N$th power descending to $C_i$; divide that degree by $N$. In particular
\begin{equation}\label{eq:root-degrees}
 \deg\bigl(\OO(\mathfrak C_j)|_{\mathfrak C_i}\bigr)=B_{ij}/N,
 \qquad \deg(J|_{\mathfrak C_i})=d_i.
\end{equation}
The degree of a vector bundle means the degree of its determinant. Frobenius multiplies it by $p$.

The restriction of $V$ to each component is semistable of degree zero. Indeed, on a finite \'etale trivializing cover the determinant of a subbundle maps into an exterior power of a trivial bundle, so has nonpositive degree. Degrees under a finite representable cover scale by a positive factor; their signs are preserved. The determinant of $V$ becomes trivial there, so has degree zero. This proof works even if the degree of the cover is divisible by $p$. The related finite-bundle statement and the root-stack interpretation of parabolic degree are treated by Borne \cite[Theorems 4--5 and Proposition 6, author pp.~12,17,22]{Borne}.

We record a cohomological consequence of this degree convention. Decompose a bundle on $\mathfrak C_i$ into characters of its generic $\mu_N$. Nontrivial characters have zero coherent cohomology, since taking invariants along that gerbe is exact. A character-zero summand $W$ descends to the rooted curve. Its coarse pushforward $E$ is a vector bundle of the same rank on $C_i$, and
\begin{equation}\label{eq:coarse-degree}
 0\leq\deg W-\deg E\leq\nu_i\operatorname{rank}W.
\end{equation}
To verify this locally, lift a character basis on a root chart by equivariant Nakayama. The adjunction map from the pullback of the invariant module has diagonal entries $z^{a_j}$, with $0\leq a_j<N$. The determinant loses $\sum_j a_j/N$ in degree. Summing at the crossing points proves \eqref{eq:coarse-degree}.

If $W$ is semistable of negative slope $t$, it has no sections, so Riemann--Roch gives
\begin{equation}\label{eq:curve-bound}
 h^1(\mathfrak C_i,W)\leq
 \operatorname{rank}W\,(g_i-1-t+\nu_i).
\end{equation}
In slope zero, descend the character-zero summand $W$ to the rooted
curve and choose a point away from its crossing points. Evaluation
there is injective: a section vanishing there would yield, after
saturation, a positive-degree line subbundle on that rooted curve. Thus in that case
\begin{equation}\label{eq:curve-zero}
 h^1(\mathfrak C_i,W)\leq
 \operatorname{rank}W\,(g_i+\nu_i).
\end{equation}
Every isotypic summand of $V$ is semistable of degree zero. Each is a direct summand and has nonpositive degree, and their degrees sum to zero. The preceding estimates therefore apply after twisting such a summand by any line bundle which cancels its generic character.

\subsection{Frobenius and a common pole cutoff}

Put $\mathcal B=F_*\OO_{\mathfrak X}/\OO_{\mathfrak X}$. This is locally free. A coefficient field makes $A$ a quotient of a power-series ring over perfect $k$, so $A$ and the finite-type chart rings are $F$-finite. Frobenius on the root stack is representable because it induces the automorphism $\zeta\mapsto\zeta^p$ on its stabilizers. It is finite flat on the regular charts, and its unit is a unimodular vector locally in $F_*\OO$.

The pullback $F^*\mathcal B$ identifies with the augmentation ideal in $F^*F_*\OO$. Its powers have locally free graded pieces the truncated symmetric powers
\[
 T^a(\Omega^1_{\mathfrak X/k}),\qquad 1\leq a\leq2(p-1).
\]
This is Sun's Frobenius filtration \cite[Theorem 3.7, author pp.~12--13]{Sun}. Its applicability here can be checked without a finite-type-over-$k$ assumption. At every closed point of the exceptional divisor, the completed regular chart is $k[[u,v]]$, and
\[
 k[[u,v]]\otimes_{k[[u^p,v^p]]}k[[u,v]]
       =k[[u,v]][\alpha,\beta]/(\alpha^p,\beta^p).
\]
The augmentation ideal is $(\alpha,\beta)$; its graded generators are $du,dv$. This parameter-independent calculation verifies the asserted finite-module filtration after faithfully flat completion, and hence near the exceptional divisor. A nonzero coherent error there would have a closed point in its support.

The conormal sequence on $\mathfrak C_i$ has line terms of degrees
\[
 -B_{ii}/N,\qquad 2g_i-2+\nu_i(1-1/N).
\]
For the second degree use $dx=Nz^{N-1}dz$ at each root point; the generic gerbe is \'etale and adds no differential direction. Each degree is at most $K_i$. The truncated symmetric powers have a filtration by monomials in this two-line flag, with both exponents at most $p-1$. Every line layer therefore has degree at most $2(p-1)K_i$. Tensoring with $V$ and using $F^*V\simeq V$ bounds the slope of every subbundle of $F^*(\mathcal B\otimes V)|_{\mathfrak C_i}$ by that same number. Pull back a subbundle by Frobenius and divide degrees by $p$. We obtain
\begin{equation}\label{eq:frobenius-slope}
 \mu(\text{any subbundle of }(\mathcal B\otimes V)|_{\mathfrak C_i})\leq b_i.
\end{equation}
The bound is independent of $N$ and $r$.

Let $\mathfrak D=\bigcup_i\mathfrak C_i$ be reduced, let $i:\mathfrak D\to\mathfrak X$, and let $j:U\to\mathfrak X$. Write $R^1\Psi M=i^*R^1j_*M$. Put $\mathfrak Z=\pi^*Z=\sum_iNc_i\mathfrak C_i$, whose ideal is again denoted $J$. Sheaves supported on a nilpotent thickening of $\mathfrak D$ are
viewed on its small \'etale site, which is invariant under such
thickenings. The inverse image $i^*$ is that of \'etale sheaves;
it preserves the infinitesimal layers of the pole quotients.
Tensor restriction of coherent modules to the reduced divisor
would discard those layers. The Artin--Schreier sequence gives
a pole filtration $U_m$ of $R^1\Psi M$, with
\begin{equation}\label{eq:pole-sequence}
 0\longrightarrow J^{-m}V/V\xrightarrow{F_V-1}J^{-pm}V/V
       \longrightarrow U_m\longrightarrow0.
\end{equation}
Here $F_V$ is the Frobenius-semilinear map given by the unit structure. We justify the exactness, including its bound on preimages. First $j_*M=M$: on a regular strict-local root chart the complement of the boundary is connected, and the extended local system is constant. On an affine chart $U$ is the complement of a principal divisor, so higher direct images of $V$ as a quasi-coherent sheaf vanish. Artin--Schreier computes $R^1\Psi M$ as the quotient of $j_*V/V$ by $F_V-1$. A positive pole of a vector, measured by its smallest component valuation in an \'etale trivialization, is multiplied by $p$ under $F_V$ and cannot cancel with its smaller pole under the identity. Thus a preimage of a vector with poles bounded by $pmZ$ has poles bounded by $mZ$. This proves \eqref{eq:pole-sequence}; it also proves that the increasing union is all of $R^1\Psi M$.

For $m\geq2$, the identity term lies in the preceding pole level because $m\leq p(m-1)$. Therefore the graded quotient is the cokernel of Frobenius between
\[
 J^{-m}V/J^{-m+1}V
       \quad\text{and}\quad J^{-pm}V/J^{-pm+p}V.
\]
Give the target its Frobenius scalar action, $a\cdot v=a^pv$. This is the coherent structure used for the following slope estimates; the underlying additive cokernel is unchanged. Since $F^{-1}\mathfrak Z=p\mathfrak Z$, projection formula and the locally free cokernel $\mathcal B$ identify it as
\begin{equation}\label{eq:pole-graded}
 U_m/U_{m-1}=
 (\mathcal B\otimes V)|_{\mathfrak Z}\otimes J^{-m}|_{\mathfrak Z}.
\end{equation}
This also follows by restricting $0\to\OO\to F_*\OO\to\mathcal B\to0$, which remains exact because $\mathcal B$ is locally free. The calculation is compatible on charts, so descends. It is the twisted root-stack form of Suzuki's calculation \cite[Proposition 3.1, author pp.~6--7]{Suzuki}.

Filter $\OO_{\mathfrak Z}$ by removing one component at a time. A line layer on $\mathfrak C_i$ has degree at most $L_i$: all coefficients of the removed cycle lie between zero and $Nc_j$, and \eqref{eq:root-degrees} applies. The corresponding layer of \eqref{eq:pole-graded} has maximal slope at most $b_i+L_i-md_i$, by \eqref{eq:frobenius-slope}. This is negative for $m\geq m_0$, so
\[
 H^0(\mathfrak D,U_m/U_{m-1})=0\qquad(m\geq m_0).
\]
Global sections commute with this filtered union on the quasi-compact exceptional stack; this can be checked on a finite affine atlas and its overlaps. It follows that
\begin{equation}\label{eq:cutoff}
 H^0(\mathfrak D,R^1\Psi M)=H^0(\mathfrak D,U_{m_0}).
\end{equation}

\subsection{Counting the low-pole contribution}

We now bound the fixed pole level without introducing a factor of $N$. We use the negative-cycle filtration underlying Suzuki's Proposition 2.1 \cite[author pp.~4--5]{Suzuki}. Its elementary argument applies to the normalized matrix $B/N$: if an effective nonzero cycle $E$ is supported on these curves, negative definiteness gives a component of its support with $E\cdot\mathfrak C_i<0$. Removing that component gives a negative-degree line layer of $\OO(E)$ modulo the preceding cycle. More generally, when the initial cycle has nonpositive intersection with every component, the same argument for the difference of cycles gives negative layers throughout. By induction this supplies filtrations of
\[
 J^{-pm}/\OO,\qquad J^{-pm}/J^{-m},\qquad J^{-m}/\OO
\]
by negative-degree line bundles on the components. Tensoring with the degree-zero semistable $V$ gives no global sections of these sheaves.

Consequently \eqref{eq:pole-sequence} injects $H^0(U_m)$ into the kernel of a Frobenius-semilinear map minus an injective linear map from $H^1(J^{-m}V/V)$ to $H^1(J^{-pm}V/V)$. If $f:W\to W'$ is Frobenius-semilinear and $g:W\to W'$ is injective and linear, this kernel has $\F_p$-dimension at most $\dim_kW$. Split $g$ linearly and reduce to fixed vectors of a semilinear endomorphism of $W$. Any minimal $k$-linear relation between fixed vectors, normalized to have one coefficient equal to one, can be subtracted from its Frobenius transform. Minimality forces every coefficient to lie in $\F_p$. Thus $\F_p$-independent fixed vectors are $k$-independent.

In a negative filtration of $J^{-m}/\OO$, a line layer $P$ on $\mathfrak C_i$ satisfies
\[
 -mL_i\leq\deg P<0.
\]
As the coefficient of $\mathfrak C_i$ runs through $mNc_i$ removals, the generic character of $P$ runs through every residue modulo $N$ exactly $mc_i$ times. Twists from other components have trivial character at its generic point. Decompose $V|_{\mathfrak C_i}=\bigoplus_\chi V_{i,\chi}$. Only the summand whose character cancels that of $P$ contributes to cohomology. The sum of its ranks over all these layers is exactly
\[
 mc_i\sum_\chi\operatorname{rank}V_{i,\chi}=mc_i r.
\]
This identity cancels the growing number of layers. Applying \eqref{eq:curve-bound} to each contributing summand, then using the long exact sequences of the filtration, gives
\begin{equation}\label{eq:low-pole-bound}
 h^1(J^{-m}V/V)\leq r m\sum_i c_i(g_i+\nu_i+mL_i+1).
\end{equation}
In particular \eqref{eq:cutoff} and this estimate with $m=m_0$ give a bound linear in $r$ for the nearby-cycle contribution.

The normalization sequence of the reduced exceptional stack has one node term for each crossing, with stabilizer $\mu_N^2$. Its invariant fibre has dimension at most $r$, and averaging kills its higher cohomology. On each normalization component only the generic-character-zero summand contributes. Equation~\eqref{eq:curve-zero} therefore gives
\begin{equation}\label{eq:exceptional-bound}
 h^1(\mathfrak D,V)\leq
 r\left(\sum_i(g_i+\nu_i)+\#\{\text{crossings}\}\right).
\end{equation}
On the finite-dimensional space $H^0(\mathfrak D,V)$, the map $F_V-1$ is surjective. Indeed its differential is minus the identity, so it is an \'etale homomorphism of additive algebraic groups with open image; an open subgroup of the connected additive group is the whole group. Artin--Schreier and the fixed-vector argument above bound $\dim_{\F_p}H^1(\mathfrak D,M)$ by \eqref{eq:exceptional-bound}.

For clarity, the nearby-cycle comparison used here follows from ordinary scheme proper base change, which applies to all torsion coefficients, including $\F_p$ \cite[Theorem 59.91.11, Tag 095S]{Stacks}. Coarse pushforward for the root stack is exact on coherent sheaves and on $\F_p$-sheaves. Locally it is invariants under a finite group of order prime to $p$, and averaging is exact. For the \'etale assertion, a strict henselian coarse neighborhood has a finite root cover which is a union of strictly henselian local schemes; their global-section functors are exact, and finite-group cohomology vanishes by averaging. This description commutes with passage to the closed fibre. Thus proper base change for $X$ over the strictly henselian $\Spec A$ gives the degree-one exact sequence
\[
 0\longrightarrow H^1(\mathfrak D,M)\longrightarrow H^1(U,M)
      \longrightarrow H^0(\mathfrak D,R^1\Psi M).
\]
Combining \eqref{eq:low-pole-bound} and \eqref{eq:exceptional-bound} proves Theorem~\ref{thm:local-coefficients}, for example with
\[
 C_X=m_0\sum_i c_i(g_i+\nu_i+m_0L_i+1)
          +\sum_i(g_i+\nu_i)+\#\{\text{crossings}\}.
\]

\section{Bounding conductors and abelian inertia}\label{sec:abelian-inertia}

We prove Theorem~\ref{thm:abelian-inertia} by turning the refined Swan conductor into an intersection inequality. Only the local integrality theorem is needed, in the published form of Kato's Theorem 7.1 restated by Bright and Newton \cite[Theorem 4.1, \S4.3, pp.~848--849]{BN}.

Fix a character $\chi\in H^1(U,\mathbf Q_p/\mathbf Z_p)$ and let $S=\sum_i s_i C_i$ be its logarithmic Swan divisor. At a closed point of the exceptional divisor the local ring of $X$ is excellent regular, its residue field is $k$, and its boundary equations form part of a regular parameter system. These are the hypotheses of the cited theorem: excellence and finite $p$-degree of the residue field suffice. They hold because $A$ is complete noetherian and $X$ is of finite type over $A$; $k$ is perfect. No finite-type-over-a-field globalization is required.

If $s_i>0$, the intrinsic generic refined conductor is nonzero. Kato's local integrality assertion at every closed point of $C_i$ extends it to a nonzero section
\begin{equation}\label{eq:rsw-section}
 0\ne\operatorname{rsw}_i(\chi)\in
 H^0\left(C_i,\Omega_X^1(\log D)|_{C_i}\otimes\OO_X(S)|_{C_i}\right).
\end{equation}
Indeed the local assertion bounds its poles along every other boundary equation by the corresponding Swan order. These are exactly the local lattices of the displayed bundle, and they agree on overlaps because the generic differential is intrinsic. The logarithmic differential bundle is locally free near $D$: this follows directly after completion to the regular two-variable power-series ring over perfect $k$.

The residue sequence on $C_i$ has the form
\[
 0\longrightarrow\Omega^1_{C_i/k}(\log\{\text{crossings}\})
 \longrightarrow\Omega_X^1(\log D)|_{C_i}
 \longrightarrow\OO_{C_i}\longrightarrow0.
\]
Its line terms have degrees $2g_i-2+\nu_i$ and zero. Hence every line subbundle of its middle term has degree at most $k_i=\max(0,2g_i-2+\nu_i)$. Saturating the map $\OO_{C_i}(-S)\to\Omega_X^1(\log D)|_{C_i}$ from \eqref{eq:rsw-section} gives $S\cdot C_i\geq-k_i$. If $s_i=0$, the same inequality holds because all off-diagonal intersections and all $s_j$ are nonnegative. In vector notation,
\[
 Bs\geq-k.
\]

For the matrix step put $C=-B$. It is positive definite with nonpositive off-diagonal entries, and $C^{-1}$ is entrywise nonnegative. To prove this, take $Cx=v\geq0$ and write $x=x_+-x_-$ with disjointly supported nonnegative parts. If $x_-\ne0$, then
\[
 0\leq x_-^Tv=x_-^TCx_+-x_-^TCx_-<0,
\]
a contradiction. Thus $x\geq0$. Applying the inverse to $Cs\leq k$ proves $s\leq(-B)^{-1}k=t$.

Choose $e$ with $p^e>\max_i t_i$. Multiplication of a character by $p$ divides its Swan bound by $p$, rounding down. Here is the characteristic-$p$ calculation. The conductor filtration is the image, under Artin--Schreier--Witt, of the Witt vectors $a=(a_{m-1},\ldots,a_0)$ satisfying $p^jv(a_j)\geq-n$ for all $j$ \cite[\S\S3.2.2--3.2.3, pp.~152--153]{KLS}. Witt restriction drops the last displayed coordinate, so every surviving weight is divided by $p$. It represents multiplication by $p$ under the inclusions into $\mathbf Q_p/\mathbf Z_p$. Therefore
\[
 s_i(p^e\chi)\leq\lfloor s_i(\chi)/p^e\rfloor=0.
\]
Conductor zero means that an integral Witt vector represents the class; its Artin--Schreier--Witt equation is \'etale over the discrete valuation ring. Thus $p^e\chi$ is unramified along $D$ and extends over $X$ by regular purity.

For a character of order $n$ prime to $p$, use the Kummer description as a line bundle on $U$ with a trivialization of its $n$th power. Extend the line bundle to the regular $X$. Its trivialization has divisor $\sum_i a_iC_i$; intersecting with every $C_j$ gives $Ba\equiv0\pmod n$. Multiplication by the adjugate of $B$ shows that $b=|\det B|$ kills all the residues $a_i$ modulo $n$, which are its tame inertia characters. Purity extends the resulting character to $X$.

The map $Q=\pi_1(U)^{\mathrm{ab}}\to\pi_1(X)^{\mathrm{ab}}$ is surjective: a connected finite \'etale cover of normal $X$ remains connected over its dense open $U$. Let $I$ be its kernel. The two character calculations imply $bp^eI=0$, since finite characters of the profinite abelian group $Q$ separate points. Suzuki's Theorem 1.2 makes $Q$ topologically finitely generated \cite[author p.~2]{Suzuki}. The bounded-exponent subgroup of a finitely generated profinite abelian group is finite. Thus $I$ is finite, proving Theorem~\ref{thm:abelian-inertia}.

\section{The maximal metabelian quotient}\label{sec:metabelian}

We deduce Theorem~\ref{thm:metabelian} from a uniform first-cohomology bound for modules with abelian monodromy. A single finite cover removes the abelian inertia for all such modules. We first give the tame generator bound used for coefficients of characteristic different from $p$.

\begin{lemma}\label{lem:tame-generators}
Let $\Delta$ be the dual graph of $D$. The full tame fundamental group relative to $(X,D)$ is generated by at most
\[
 d_X=\sum_i(2g_i+\nu_i+1)+b_1(\Delta)
\]
elements.
\end{lemma}
\begin{proof}
Every finite tame cover extends to a root stack $\mathfrak X$ of some prime-to-$p$ order $N$, as above. The reduced exceptional stack $\mathfrak D$ surjects on its fundamental group. To see this, a connected finite \'etale cover of $\mathfrak X$ has coarse space finite over $X$: on a root chart its invariant ring is a submodule of a finite module over a noetherian ring. This coarse space is proper over the henselian local base $A$. Its closed fibre is connected by proper henselian connectedness \cite[Tag 0A48]{Stacks}. Taking invariants commutes with reduction by averaging, so that fibre is the coarse space of the closed-fibre cover of $\mathfrak X$. The latter, and hence its reduction over $\mathfrak D$, is connected. This proves the asserted surjectivity.

On a normalization component $\mathfrak C_i$, the generic $\mu_N$ is central in the monodromy of every finite cover: its band acts by natural automorphisms of the fibre functor. Quotienting this cyclic image makes the cover descend to the rooted curve. Restriction to the complement of its crossing points gives a tame cover of the ordinary curve, and a connected cover stays connected there. SGA~1, Expos\'e XIII, Corollaire 2.12 bounds this full tame curve group by $2g_i+\nu_i$ generators \cite[original pp.~392--395; recomposed pp.~290--292]{SGA1}. Adding one generator for generic inertia gives $2g_i+\nu_i+1$ for the component group.

Finite \'etale covers of $\mathfrak D$ glue from its normalization with identifications at the two preimages of each node. On a local chart this is descent of finite $\mu_N^2$-sets on the strictly henselian nodal ring $k[[u,v]]/(uv)$ and its two strictly henselian branches. Choose a spanning tree of $\Delta$. Its identifications give a single fibre; the remaining edges each contribute one gluing element. The images of the component groups and these elements generate every connected finite monodromy group. Otherwise their orbits would decompose all the component covers compatibly with the gluing, contradicting connectedness. This proves the displayed bound for every finite quotient, independently of $N$. Compactness of the spaces of generating tuples proves the same bound for the profinite group.
\end{proof}

\begin{lemma}\label{lem:abelian-coefficients}
There is a constant $C_A$ such that, for every prime $\ell$ and every finite-dimensional $\F_\ell$-local system $M$ on $U$ with abelian monodromy,
\[
 \dim_{\F_\ell}H^1(U,M)\leq C_A\dim_{\F_\ell}M.
\]
\end{lemma}
\begin{proof}
The finite group $I$ of Theorem~\ref{thm:abelian-inertia} is separated from zero by a finite quotient of $Q=\pi_1(U)^{\mathrm{ab}}$. Thus choose an open subgroup $H\subset Q$ with $H\cap I=0$, and let $U'\to U$ be the associated fixed connected finite Galois cover, of group $G$. The integral closure $A'$ of $A$ in its function field is finite, complete, normal and local, with residue field $k$. Choose an SNC resolution $X'$ of $\Spec A'$ dominating the normalization of $X$ in that field. Exceptional inertia on $X'$ maps into $I$, because every \'etale cover of $X$ pulls back unramified there; it also maps into $H$, by the choice of $U'$. Its image in $Q$ is therefore zero.

Every abelian-monodromy $M$ factors through $Q$, so its pullback $M'$ to $U'$ is unramified along the boundary and extends to $X'$ by purity. When $\ell=p$, Theorem~\ref{thm:local-coefficients} on $X'$ bounds its first cohomology by a fixed multiple of its rank. When $\ell\ne p$, its first-cohomology classes factor through the full tame quotient. Indeed a continuous cocycle defines a finite extension of its monodromy by the additive $\F_\ell$-module. On wild inertia the original action is trivial and the cocycle is a homomorphism from a pro-$p$ group to an $\ell$-group, hence zero. Lemma~\ref{lem:tame-generators} now bounds the cocycle space by $d_{X'}\operatorname{rank}M'$: evaluation on generators is injective.

Hochschild--Serre for $U'/U$ injects the kernel of restriction in degree one into $H^1(G,H^0(U',M'))$. This has dimension at most $d(G)\operatorname{rank}M$, again by evaluation on a generating set. Adding this fixed bound to the two bounds on $U'$ proves the assertion, uniformly in $\ell$.
\end{proof}

\begin{proof}[Proof of Theorem~\ref{thm:metabelian}]
Let $d_0$ be a generator bound for $\pi_1(U)^{\mathrm{ab}}$, and choose an integer $D\geq\max\{d_0,C_A+1\}$. We show by induction on the order that every finite metabelian quotient $F$ of $\pi_1(U)$ is generated by $D$ elements. If $F$ is abelian this follows from $d_0$. Otherwise choose a minimal nonzero normal subgroup $M$ inside its abelian derived subgroup. Then $M$ is an elementary abelian $\ell$-group for some prime, simple as an $F$-module under conjugation. Its action factors through $F^{\mathrm{ab}}$.

Choose a generating $D$-tuple in $F/M$, padding with identities if necessary. Every lift generates a subgroup $K$ with $KM=F$. The intersection $K\cap M$ is normal in $F$: it is normalized by $K$, and by $M$ because $M$ is abelian. Simplicity gives $K\cap M=M$ or zero. In the first case $K=F$; in the second it is a complement to $M$. Thus failure of all lifts could occur only if $F$ splits as $M\rtimes H$ with $H=F/M$.

In the split case,
\[
 H^1(F,M)=H^1(H,M)\oplus\operatorname{End}_H(M).
\]
Indeed restricting a cocycle to $M$ gives an $H$-equivariant endomorphism, and each such endomorphism extends by $c(mh)=c(m)$; the remaining cocycle comes from $H$. The second summand is nonzero. Inflation from $F$ to $\pi_1(U)$ is injective, so Lemma~\ref{lem:abelian-coefficients} gives
\[
 \dim_{\F_\ell}H^1(H,M)\leq C_A\dim_{\F_\ell}M-1.
\]
Complements correspond to crossed homomorphisms $H\to M$. Their number is
\[
 |Z^1(H,M)|=|H^1(H,M)|\,|M|/|M^H|
       \leq |M|^{C_A+1}/\ell < |M|^D.
\]
There are $|M|^D$ lifts of the chosen tuple, and each complement contains exactly one such lift. A generating lift therefore exists. This proves the induction. Finally the closed sets of $D$-tuples generating each finite quotient have the finite intersection property in the compact space of $D$-tuples of the maximal metabelian quotient. Their intersection is nonempty and consists of topological generating tuples.
\end{proof}
